% GUID v6, Primer P5 (label box:holo): holography and quantum error correction. Background only; it states no result of the paper.
% Place it at the start of Section sec:holo (marker %PRIMER:holo in part1.tex): % GUID v6, Primer P5 (label box:holo): holography and quantum error correction. Background only; it states no result of the paper.
% Place it at the start of Section sec:holo (marker %PRIMER:holo in part1.tex): % GUID v6, Primer P5 (label box:holo): holography and quantum error correction. Background only; it states no result of the paper.
% Place it at the start of Section sec:holo (marker %PRIMER:holo in part1.tex): % GUID v6, Primer P5 (label box:holo): holography and quantum error correction. Background only; it states no result of the paper.
% Place it at the start of Section sec:holo (marker %PRIMER:holo in part1.tex): \input{primers/p5-holo}
\begin{primer}{Holography and quantum error correction}{box:holo}
\emph{Stabiliser codes.} A quantum error-correcting code stores $k$ logical qubits in $n$ physical ones so that some errors can be
undone. In a \emph{stabiliser code} the code space is the joint $+1$ eigenspace of $n-k$ commuting products of Pauli matrices, the
\emph{stabilisers}~\cite{gottesman1997}. Measuring them gives the \emph{syndrome}. An error that anticommutes with a stabiliser flips its
sign and so shows in the syndrome. The stabilisers commute with the logical operators, so the syndrome reveals nothing about the stored
content. A code of distance $d$, written $[[n,k,d]]$, corrects any $\lfloor(d-1)/2\rfloor$ unknown errors, or any $d-1$
\emph{erasures} (losses at known places).

\emph{The five-qubit code.} The smallest code that corrects every single-qubit error is $[[5,1,3]]$, with the stabiliser $XZZXI$ and its
cyclic shifts~\cite{laflamme,bennett1996}. It corrects any two erasures, so its logical qubit can be recovered from any three of the five
qubits and from no two: a $((3,5))$ quantum threshold scheme~\cite{cgl}.

\emph{Perfect tensors and the HaPPY code.} A tensor with $2m$ legs is \emph{perfect} if every split of its legs into $m$ and $m$ is, up to
normalisation, a unitary map. The encoding map of the five-qubit code, with its logical leg as a sixth, is such a tensor. Pastawski,
Yoshida, Harlow and Preskill tile the hyperbolic disc with pentagons, place one perfect tensor on each, contract neighbouring legs, and
leave one leg per tile in the bulk and the outermost legs on the boundary~\cite{happy}. The network is an isometry from the bulk to the
boundary, a toy model of the bulk-boundary correspondence of AdS/CFT~\cite{maldacena}.

\emph{Ryu--Takayanagi as a minimal cut.} In AdS/CFT the entropy of a boundary region $A$ is $S(A)=\mathrm{Area}(\gamma_A)/4G\hbar$ (with
$c=k_B=1$), where $\gamma_A$ is the bulk surface of least area anchored on the edge of $A$~\cite{rt}. In a network of perfect tensors the
entropy of $A$, in bits, is for many regions the number of qubit legs cut by the smallest cut that separates $A$ from the rest of the
boundary~\cite{happy}.

\emph{Reconstruction.} The bulk region between $A$ and $\gamma_A$ is the \emph{entanglement wedge} of $A$. Every bulk operator inside it can
be represented by an operator acting on $A$ alone~\cite{adh,dong2016}: the bulk is encoded redundantly, like a logical qubit. A code has
\emph{complementary recovery} when what cannot be reconstructed on $A$ can be reconstructed on the rest of the boundary.

\emph{The area operator.} Harlow showed that every code with complementary recovery obeys a quantum Ryu--Takayanagi formula
(Eq.~\eqref{eq:harlow} in Section~\ref{sec:holo}). Its area term is the expectation value of an \emph{area operator} that lies in the centre
of the algebra reconstructable on $A$: it commutes with everything that either side can do~\cite{harlow}. A \emph{fixed-area state} is an
eigenstate of this operator, and such states have a flat entanglement spectrum~\cite{dhm}. A general state of quantum gravity is a
superposition of areas.
\end{primer}

\begin{primer}{Holography and quantum error correction}{box:holo}
\emph{Stabiliser codes.} A quantum error-correcting code stores $k$ logical qubits in $n$ physical ones so that some errors can be
undone. In a \emph{stabiliser code} the code space is the joint $+1$ eigenspace of $n-k$ commuting products of Pauli matrices, the
\emph{stabilisers}~\cite{gottesman1997}. Measuring them gives the \emph{syndrome}. An error that anticommutes with a stabiliser flips its
sign and so shows in the syndrome. The stabilisers commute with the logical operators, so the syndrome reveals nothing about the stored
content. A code of distance $d$, written $[[n,k,d]]$, corrects any $\lfloor(d-1)/2\rfloor$ unknown errors, or any $d-1$
\emph{erasures} (losses at known places).

\emph{The five-qubit code.} The smallest code that corrects every single-qubit error is $[[5,1,3]]$, with the stabiliser $XZZXI$ and its
cyclic shifts~\cite{laflamme,bennett1996}. It corrects any two erasures, so its logical qubit can be recovered from any three of the five
qubits and from no two: a $((3,5))$ quantum threshold scheme~\cite{cgl}.

\emph{Perfect tensors and the HaPPY code.} A tensor with $2m$ legs is \emph{perfect} if every split of its legs into $m$ and $m$ is, up to
normalisation, a unitary map. The encoding map of the five-qubit code, with its logical leg as a sixth, is such a tensor. Pastawski,
Yoshida, Harlow and Preskill tile the hyperbolic disc with pentagons, place one perfect tensor on each, contract neighbouring legs, and
leave one leg per tile in the bulk and the outermost legs on the boundary~\cite{happy}. The network is an isometry from the bulk to the
boundary, a toy model of the bulk-boundary correspondence of AdS/CFT~\cite{maldacena}.

\emph{Ryu--Takayanagi as a minimal cut.} In AdS/CFT the entropy of a boundary region $A$ is $S(A)=\mathrm{Area}(\gamma_A)/4G\hbar$ (with
$c=k_B=1$), where $\gamma_A$ is the bulk surface of least area anchored on the edge of $A$~\cite{rt}. In a network of perfect tensors the
entropy of $A$, in bits, is for many regions the number of qubit legs cut by the smallest cut that separates $A$ from the rest of the
boundary~\cite{happy}.

\emph{Reconstruction.} The bulk region between $A$ and $\gamma_A$ is the \emph{entanglement wedge} of $A$. Every bulk operator inside it can
be represented by an operator acting on $A$ alone~\cite{adh,dong2016}: the bulk is encoded redundantly, like a logical qubit. A code has
\emph{complementary recovery} when what cannot be reconstructed on $A$ can be reconstructed on the rest of the boundary.

\emph{The area operator.} Harlow showed that every code with complementary recovery obeys a quantum Ryu--Takayanagi formula
(Eq.~\eqref{eq:harlow} in Section~\ref{sec:holo}). Its area term is the expectation value of an \emph{area operator} that lies in the centre
of the algebra reconstructable on $A$: it commutes with everything that either side can do~\cite{harlow}. A \emph{fixed-area state} is an
eigenstate of this operator, and such states have a flat entanglement spectrum~\cite{dhm}. A general state of quantum gravity is a
superposition of areas.
\end{primer}

\begin{primer}{Holography and quantum error correction}{box:holo}
\emph{Stabiliser codes.} A quantum error-correcting code stores $k$ logical qubits in $n$ physical ones so that some errors can be
undone. In a \emph{stabiliser code} the code space is the joint $+1$ eigenspace of $n-k$ commuting products of Pauli matrices, the
\emph{stabilisers}~\cite{gottesman1997}. Measuring them gives the \emph{syndrome}. An error that anticommutes with a stabiliser flips its
sign and so shows in the syndrome. The stabilisers commute with the logical operators, so the syndrome reveals nothing about the stored
content. A code of distance $d$, written $[[n,k,d]]$, corrects any $\lfloor(d-1)/2\rfloor$ unknown errors, or any $d-1$
\emph{erasures} (losses at known places).

\emph{The five-qubit code.} The smallest code that corrects every single-qubit error is $[[5,1,3]]$, with the stabiliser $XZZXI$ and its
cyclic shifts~\cite{laflamme,bennett1996}. It corrects any two erasures, so its logical qubit can be recovered from any three of the five
qubits and from no two: a $((3,5))$ quantum threshold scheme~\cite{cgl}.

\emph{Perfect tensors and the HaPPY code.} A tensor with $2m$ legs is \emph{perfect} if every split of its legs into $m$ and $m$ is, up to
normalisation, a unitary map. The encoding map of the five-qubit code, with its logical leg as a sixth, is such a tensor. Pastawski,
Yoshida, Harlow and Preskill tile the hyperbolic disc with pentagons, place one perfect tensor on each, contract neighbouring legs, and
leave one leg per tile in the bulk and the outermost legs on the boundary~\cite{happy}. The network is an isometry from the bulk to the
boundary, a toy model of the bulk-boundary correspondence of AdS/CFT~\cite{maldacena}.

\emph{Ryu--Takayanagi as a minimal cut.} In AdS/CFT the entropy of a boundary region $A$ is $S(A)=\mathrm{Area}(\gamma_A)/4G\hbar$ (with
$c=k_B=1$), where $\gamma_A$ is the bulk surface of least area anchored on the edge of $A$~\cite{rt}. In a network of perfect tensors the
entropy of $A$, in bits, is for many regions the number of qubit legs cut by the smallest cut that separates $A$ from the rest of the
boundary~\cite{happy}.

\emph{Reconstruction.} The bulk region between $A$ and $\gamma_A$ is the \emph{entanglement wedge} of $A$. Every bulk operator inside it can
be represented by an operator acting on $A$ alone~\cite{adh,dong2016}: the bulk is encoded redundantly, like a logical qubit. A code has
\emph{complementary recovery} when what cannot be reconstructed on $A$ can be reconstructed on the rest of the boundary.

\emph{The area operator.} Harlow showed that every code with complementary recovery obeys a quantum Ryu--Takayanagi formula
(Eq.~\eqref{eq:harlow} in Section~\ref{sec:holo}). Its area term is the expectation value of an \emph{area operator} that lies in the centre
of the algebra reconstructable on $A$: it commutes with everything that either side can do~\cite{harlow}. A \emph{fixed-area state} is an
eigenstate of this operator, and such states have a flat entanglement spectrum~\cite{dhm}. A general state of quantum gravity is a
superposition of areas.
\end{primer}

\begin{primer}{Holography and quantum error correction}{box:holo}
\emph{Stabiliser codes.} A quantum error-correcting code stores $k$ logical qubits in $n$ physical ones so that some errors can be
undone. In a \emph{stabiliser code} the code space is the joint $+1$ eigenspace of $n-k$ commuting products of Pauli matrices, the
\emph{stabilisers}~\cite{gottesman1997}. Measuring them gives the \emph{syndrome}. An error that anticommutes with a stabiliser flips its
sign and so shows in the syndrome. The stabilisers commute with the logical operators, so the syndrome reveals nothing about the stored
content. A code of distance $d$, written $[[n,k,d]]$, corrects any $\lfloor(d-1)/2\rfloor$ unknown errors, or any $d-1$
\emph{erasures} (losses at known places).

\emph{The five-qubit code.} The smallest code that corrects every single-qubit error is $[[5,1,3]]$, with the stabiliser $XZZXI$ and its
cyclic shifts~\cite{laflamme,bennett1996}. It corrects any two erasures, so its logical qubit can be recovered from any three of the five
qubits and from no two: a $((3,5))$ quantum threshold scheme~\cite{cgl}.

\emph{Perfect tensors and the HaPPY code.} A tensor with $2m$ legs is \emph{perfect} if every split of its legs into $m$ and $m$ is, up to
normalisation, a unitary map. The encoding map of the five-qubit code, with its logical leg as a sixth, is such a tensor. Pastawski,
Yoshida, Harlow and Preskill tile the hyperbolic disc with pentagons, place one perfect tensor on each, contract neighbouring legs, and
leave one leg per tile in the bulk and the outermost legs on the boundary~\cite{happy}. The network is an isometry from the bulk to the
boundary, a toy model of the bulk-boundary correspondence of AdS/CFT~\cite{maldacena}.

\emph{Ryu--Takayanagi as a minimal cut.} In AdS/CFT the entropy of a boundary region $A$ is $S(A)=\mathrm{Area}(\gamma_A)/4G\hbar$ (with
$c=k_B=1$), where $\gamma_A$ is the bulk surface of least area anchored on the edge of $A$~\cite{rt}. In a network of perfect tensors the
entropy of $A$, in bits, is for many regions the number of qubit legs cut by the smallest cut that separates $A$ from the rest of the
boundary~\cite{happy}.

\emph{Reconstruction.} The bulk region between $A$ and $\gamma_A$ is the \emph{entanglement wedge} of $A$. Every bulk operator inside it can
be represented by an operator acting on $A$ alone~\cite{adh,dong2016}: the bulk is encoded redundantly, like a logical qubit. A code has
\emph{complementary recovery} when what cannot be reconstructed on $A$ can be reconstructed on the rest of the boundary.

\emph{The area operator.} Harlow showed that every code with complementary recovery obeys a quantum Ryu--Takayanagi formula
(Eq.~\eqref{eq:harlow} in Section~\ref{sec:holo}). Its area term is the expectation value of an \emph{area operator} that lies in the centre
of the algebra reconstructable on $A$: it commutes with everything that either side can do~\cite{harlow}. A \emph{fixed-area state} is an
eigenstate of this operator, and such states have a flat entanglement spectrum~\cite{dhm}. A general state of quantum gravity is a
superposition of areas.
\end{primer}
